\section{Future Work}
\label{sec:futurework}

There are several natural next steps.

\myparagraph{Extend the \ir fragment.}
The most obvious direction is to add the features from the proposal that are
still missing: memory, pointers, \texttt{load}, \texttt{store},
\texttt{select}, and eventually limited control flow. Doing this well would
require revisiting both the state model and the well-formedness judgments.

\myparagraph{Implement a real optimizer layer.}
The current project proves refinement between hand-written pairs of
definitions. A next step is to define a syntax-directed peephole optimizer and
prove that each rewrite it performs is sound with respect to the existing
refinement relation. The current theorem library provides a starting point for
that.

\myparagraph{Improve the model of \poison and \vundef.}
The current \vundef supply is explicit and useful, but not fully faithful to
\llvm. Likewise, \texttt{freeze} is simplified. A stronger future project would
study a more principled semantics for these features and compare it against the
current executable approximation.

\myparagraph{Develop proof automation.}
At the moment, proofs are still mostly hand-guided. There is room for tactics
or theorem schemas that discharge routine refinement obligations automatically,
especially once more optimization patterns are added.

\section{Conclusion}
\label{sec:conclusion}

\proj does not yet verify a realistic fragment of all
of \llvm, and it does not implement the full optimizer-oriented roadmap set out
in the proposal. What it does provide is a compact, readable, executable \lean
formalization in which small \llvm-like programs can be written directly,
instantiated at arbitrary widths, executed, and related by machine-checked
refinement proofs. In particular, the project now supports examples involving
symbolic widths and explicit \vundef behavior, which turned out to be a richer
and more educational result than a purely arithmetic toy model would have
been.
